\documentclass[10pt,a4paper]{amsart}
\usepackage[margin=19mm]{geometry}
\usepackage{amsmath,amssymb,mathtools}
\usepackage[scr=rsfs,cal=euler]{mathalfa}
\usepackage{xfrac}
\usepackage[unicode,hidelinks,bookmarks=false]{hyperref}
\newcommand{\ie}{i.e.\ }
\newcommand{\cf}{cf.\ }
\newcommand{\resp}{respectively\ }
\newcommand{\Subsection}{\subsection}
\providecommand{\nobreakdash}{\nobreak\hbox{-}}
\setlength{\emergencystretch}{2em}
\multlinegap0pt
\newcommand{\chg}[1]{#1}
\usepackage[shortlabels]{enumitem}
\usepackage{amssymb}
\usepackage{mathtools}
\usepackage{xcolor}
\usepackage{bm}
\newtheorem{theorem}{Theorem}
\newtheorem{proposition}{Proposition}
\newtheorem{lemma}{Lemma}
\newcommand{\Id}{\mathrm{Id}}
\newcommand{\abs}[1]{\left|#1\right|}
\newcommand{\angles}[1]{\left\langle #1 \right\rangle}
\newcommand{\bc}{\mathrm{bc}}
\newcommand{\cE}{\mathcal E}
\newcommand{\cG}{\mathcal G}
\newcommand{\dd}{\mathrm d}
\newcommand{\diag}{\operatorname{diag}}
\newcommand{\disc}{\operatorname{disc}}
\newcommand{\ee}{\mathrm e}
\newcommand{\eps}{\varepsilon}
\newcommand{\ii}{\mathrm i}
\newcommand{\kap}{\kappa}
\newcommand{\la}{\lambda}
\newcommand{\om}{\omega}
\newcommand{\vp}{\varpi}
\title{\textbf{\chg{Threshold asymptotics and decay for massive Maxwell on subextremal Reissner--Nordstr\"om}}}
\author{Bobby Eka Gunara}
\date{Source: arXiv:2603.23905v1 (CC BY 4.0)}
\begin{document}
\maketitle

\noindent\emph{Source and licence.} \textbf{\chg{Threshold asymptotics and decay for massive Maxwell on subextremal Reissner--Nordstr\"om}}. Version arXiv:2603.23905v1 (CC BY 4.0), distributed by arXiv under the Creative Commons Attribution 4.0 International licence, http://creativecommons.org/licenses/by/4.0/, as recorded in the arXiv licence metadata for this exact version and shown on the abstract page https://arxiv.org/abs/2603.23905v1. Full licence notice: \texttt{licenses/arxiv-2603.23905-CC-BY-4.0.txt}.\par\smallskip

\noindent\textit{Editorial scope.} Counted unit: the article's theorem thm:uniform\_angular (source0.tex lines 437--458). The article's complete original proof of it is delivered with it (source0.tex lines 2285--2426, one \texttt{proof} environment of 142 lines, which the article places away from the statement). No mathematics is written, repaired or re-derived: the statement and the proof are the article's own lines, in the article's own notation, and no retained line is substituted or re-flowed.  Retained uncounted as the source-local context the proof needs: lines 502--513; lines 981--989; lines 1033--1064; lines 1295--1323; lines 1776--1792; lines 1798--1810; lines 1853--1865; lines 1926--1926; lines 2012--2024; lines 2030--2043; lines 2092--2111; lines 2121--2140; lines 2182--2212; lines 2222--2239; lines 2245--2254; lines 2260--2270.  Nothing was omitted from any retained span, so the package carries no omission record.  No retained line contains a citation, so the wrapper carries no bibliography and no bibliography entry.  No retained line contains a figure environment, an image-inclusion command or drawing code, so no image is delivered and the package declares no asset dependency.  Editorial formatting only, in addition: an AMS \texttt{amsart} preamble that restates the article's own declarations and macros used by the retained lines, namely the environments \texttt{theorem}, \texttt{proposition}, \texttt{lemma} on separate counters, together with the standard packages those lines need.  The retained environments print in this document as Theorem 1, Proposition 1, Lemma 1 in order of appearance, whereas the article prints the same items with the numbering of its own full document.  The complete native source of the article as distributed in the arXiv e-print for this exact version is bundled unchanged as \texttt{sources/197119/source0.tex}.
\begin{theorem}[Mode stability and threshold exclusion]\label{thm:mode_stability}
Fix $\ell\ge1$.  There is no nontrivial separated solution of the form
\[
   A=\ee^{-\ii\om t}\widehat A(r,\omega)
\]
which is ingoing at the future event horizon and outgoing or decaying at infinity, with frequency $\om$ in the closed upper half-plane.  More precisely:
\begin{enumerate}[label=\textnormal{(\roman*)}]
   \item if $\Im\om>0$, no such mode exists;
   \item if $\om\in[-\mu,\mu]$, there is no real cut pole and no threshold resonance at $\om=\pm\mu$.
\end{enumerate}
Consequently the Evans determinant $\cE_\ell$ is zero-free on $\{\Im\om>0\}$ and zero-free in punctured threshold neighborhoods of $\om=\pm\mu$.
\end{theorem}

\begin{align}
   u_2''+\Bigl(\om^2-f\mu^2-f\frac{\la^2}{r^2}\Bigr)u_2
   -\frac{2f}{r^2}\Bigl(f-\frac{r f'}{2}\Bigr)(u_2-u_3)&=0,
   \label{eq:even_frequency_domain_1}
   \\
   u_3''+\Bigl(\om^2-f\mu^2-f\frac{\la^2}{r^2}\Bigr)u_3
   +\frac{2f\la^2}{r^2}u_2&=0,
   \label{eq:even_frequency_domain_2}
\end{align}

\begin{proposition}[Reissner--Nordstr\"om even-sector polarization form]\label{prop:even_polarization_form}
For every $\ell\ge1$, the even frequency-domain system
\eqref{eq:even_frequency_domain_1}--\eqref{eq:even_frequency_domain_2} takes the matrix form
\begin{equation}\label{eq:even_matrix_form}
   \bm v''+\bigl(\om^2-f\mu^2\bigr)\bm v
   -\frac{f}{r^2}D_\ell\bm v
   -\frac{6Mf}{(2\ell+1)r^3}C_\ell\bm v
   +\frac{4Q^2f}{(2\ell+1)r^4}C_\ell\bm v
   =0,
\end{equation}
where
\begin{equation}\label{eq:D_matrix}
   \bm v=\binom{v_{-1}}{v_{+1}},
   \qquad
   D_\ell=\diag\bigl(\ell(\ell-1),(\ell+1)(\ell+2)\bigr),
\end{equation}
and
\begin{equation}\label{eq:C_matrix}
   C_\ell=
   \begin{pmatrix}
      \ell^2&-\ell(\ell+1)\\
      \ell(\ell+1)&-(\ell+1)^2
   \end{pmatrix}.
\end{equation}
In particular, the leading $r^{-2}$ part diagonalizes exactly, with effective angular momenta
\begin{equation}\label{eq:effective_Ls}
   L_{-1}=\ell-1,
   \qquad
   L_{+1}=\ell+1,
\end{equation}
and the charge $Q$ first enters the even coupling at order $r^{-4}$.
\end{proposition}

\begin{lemma}[Far-zone threshold normal form]\label{lem:normal_form}
Fix $\ell\ge1$.  There exist $R_0>r_+$, $\eta>0$, a smooth invertible transformation
\[
   U(r,\om)=S(r,\om)W(r,\om),
   \qquad
   S(r,\om)=\Id+O(r^{-1}),
\]
and channel-dependent indices $\nu_{\ell,P}>-1/2$ such that for $r\ge R_0$ every odd/even polarization channel is reduced to
\begin{equation}\label{eq:threshold_normal_form}
   w_P''+
   \Bigl(
      -\vp^2
      +\frac{2M\mu^2}{r}
      -\frac{\nu_{\ell,P}^2-\tfrac14}{r^2}
      +V_{\ell,P}^{\mathrm{sr}}(r,\om)
   \Bigr)w_P
   =0,
\end{equation}
where $\vp=\sqrt{\mu^2-\om^2}$ and
\[
   \partial_\om^jV_{\ell,P}^{\mathrm{sr}}(r,\om)=O(r^{-3}),
   \qquad j=0,1,
\]
uniformly for $\abs{\Im\om}<\eta$.  Moreover,
\begin{equation}\label{eq:nu_expansion}
   \nu_{\ell,P}=L_P+\frac12+O\bigl((M\mu)^2+(Q\mu)^2\bigr)
\end{equation}
in the small-mass regime.
\end{lemma}

\begin{lemma}[Small-$\kap$ outgoing basis]\label{lem:small_kappa_basis}
Assume $\kap\le\kap_1$ and $\vp\le\vp_1$ with $\kap_1,\vp_1$ sufficiently small.  Then the outgoing basis has the representation
\begin{equation}\label{eq:small_kappa_basis_rep}
   g_{\ell,P}(r,\om)
   =
   \sqrt{\vp r}\,K_{\nu_{\ell,P}}(\vp r)\bigl(1+E_K(r,\om)\bigr)
   +
   \sqrt{\vp r}\,I_{\nu_{\ell,P}}(\vp r)\,E_I(r,\om),
\end{equation}
where
\[
   E_K(r,\om)=O(\kap)+O(\vp^2),
   \qquad
   E_I(r,\om)=O(\kap)+O(\vp^2)
\]
uniformly for $r$ in compact subsets of $(r_+,\infty)$.
\end{lemma}

\begin{proposition}[Small-$\kap$ threshold discontinuity]\label{prop:small_kappa_disc}
In the regime $\kap\ll1$ one has
\begin{equation}\label{eq:small_kappa_disc}
   \disc\,\cG_{\ell,P}(\om;r,r')
   =
   a_{\ell,P}(r,r')\,\vp^{2\nu_{\ell,P}}
   +
   O\!\bigl(\kap\,\vp^{2\nu_{\ell,P}}\bigr)
   +
   O\!\bigl(\vp^{2\nu_{\ell,P}+2}\bigr)
\end{equation}
uniformly for $r,r'$ in compact subsets of $(r_+,\infty)$.
\end{proposition}

\begin{proposition}[Large-$\kap$ threshold discontinuity]\label{prop:large_kappa_disc}
In the regime $\kap\gg1$ one has
\begin{equation}\label{eq:large_kappa_disc}
   \disc\,\cG_{\ell,P}(\om;r,r')
   =
   b^+_{\ell,P}(r,r')\,\ee^{2\pi\ii\kap}
   +
   b^-_{\ell,P}(r,r')\,\ee^{-2\pi\ii\kap}
   +
   O(\kap^{-1})+O(\vp)
\end{equation}
uniformly for $r,r'$ in compact subsets of $(r_+,\infty)$.
\end{proposition}

\section{Explicit branch-cut tails}\label{sec:threshold_to_time}

\begin{lemma}[Endpoint oscillatory integral]\label{lem:endpoint_integral}
Let $\alpha>-1$ and $a>0$.  Then, as $t\to\infty$,
\begin{equation}\label{eq:endpoint_integral}
   \int_0^\eps\ee^{\ii a\vp^2 t}\vp^\alpha\,\dd\vp
   =
   \frac12
   \Gamma\Bigl(\frac{\alpha+1}{2}\Bigr)
   \ee^{\frac{\pi\ii}{4}(\alpha+1)}
   a^{-(\alpha+1)/2}
   t^{-(\alpha+1)/2}
   +O\bigl(t^{-(\alpha+3)/2}\bigr).
\end{equation}
\end{lemma}

\begin{theorem}[Intermediate branch-cut asymptotics]\label{thm:intermediate}
If $\kap_*(t)=M\mu^{3/2}t^{1/2}\to0$, then
\begin{align}\label{eq:intermediate_explicit}
   u_{\ell m,P}^{\bc}(t,r,r')
   &=
   A_{\ell,P}(r,r';Q)\,t^{-(\nu_{\ell,P}+1)}
   \sin(\mu t+\delta_{\ell,P}(Q))
   \notag\\
   &\qquad
   +O\!\bigl(\kap_*(t)t^{-(\nu_{\ell,P}+1)}\bigr)
   +O\!\bigl(t^{-(\nu_{\ell,P}+2)}\bigr),
\end{align}
uniformly for $r,r'$ in compact subsets of $(r_+,\infty)$.
\end{theorem}

\begin{proposition}[Model saddle estimate]\label{prop:saddle_estimate}
Let
\begin{equation}\label{eq:generic_saddle_integral}
   I(t)=\int_0^\eps(\beta\vp+O(\vp^2))\ee^{\ii\Psi_t(\vp)}\,\dd\vp,
\end{equation}
where
\begin{equation}\label{eq:phase_function}
   \Psi_t(\vp)=\frac{t\vp^2}{2\mu}+\frac{2\pi M\mu^2}{\vp}.
\end{equation}
Then, as $t\to\infty$,
\begin{equation}\label{eq:saddle_estimate}
   I(t)
   =
   \beta\,\vp_0(t)\sqrt{\frac{2\pi\mu}{3t}}\,
   \ee^{\ii\Psi_t(\vp_0(t))+\pi\ii/4}
   +
   O\!\bigl((\kap_0(t)^{-1}+\vp_0(t))t^{-5/6}\bigr),
\end{equation}
where $\vp_0(t)=(2\pi M\mu^3/t)^{1/3}$.
\end{proposition}

\begin{theorem}[Universal very-late asymptotics]\label{thm:late}
If $\kap_0(t)\to\infty$, then
\begin{align}\label{eq:late_explicit_sine}
   u_{\ell m,P}^{\bc}(t,r,r')
   &=
   B_{\ell,P}(r,r';Q)\,t^{-5/6}
   \sin\!\Bigl(
      \mu t
      -\frac32(2\pi M\mu)^{2/3}(\mu t)^{1/3}
      \notag\\
   &\qquad\qquad
      +\delta_{\ell,P,0}(Q)
      +O(\kap_0(t)^{-1}+\vp_0(t))
   \Bigr)
   \notag\\
   &\qquad
   +O\!\bigl(\kap_0(t)^{-1}+\vp_0(t)\bigr)\,t^{-5/6}.
\end{align}
The exponent $5/6$ is independent of $\ell$, of the polarization, and of the black-hole charge $Q$.
\end{theorem}

\begin{lemma}[Large-$\ell$ barrier geometry and turning points]\label{lem:largeell_barrier}
Fix a compact frequency interval $J\Subset[-\mu,\mu]$ and a compact radial set $K\Subset(r_+,\infty)$.  There exist $h_0>0$, a cutoff neighborhood $\widetilde K\Subset(r_+,\infty)$ containing $K$, and smooth functions
\[
   r_{-,P}(\om;h)<r_{+,P}(\om;h),
   \qquad
   (\om,h)\in J\times(0,h_0),
   \quad P\in\{-1,0,+1\},
\]
such that the following hold uniformly in $P$:
\begin{enumerate}[label=\textnormal{(\roman*)}]
   \item the scalarized channel symbol
   \[
      q_{P}(r,\xi;\om,h)
      :=
      \xi^2+V_{0,P}(r)+hV_{1,P}(r,\om)+h^2V_{2,P}(r,\om)-h^2\om^2
   \]
   vanishes at $\xi=0$ exactly at the two turning points $r_{-,P}(\om;h)$ and $r_{+,P}(\om;h)$;
   \item the turning points are simple and satisfy
   \[
      r_{-,P}(\om;h)-r_+\asymp h^2,
      \qquad
      r_{+,P}(\om;h)\asymp h^{-1};
   \]
   \item one has
   \[
      q_P(r,\xi;\om,h)\ge c_K(1+\xi^2)
      \qquad\text{for }r\in\widetilde K,
   \]
   so no turning point intersects the compact region supporting the cut-off resolvent.
\end{enumerate}
\end{lemma}

\begin{proposition}[Uniform WKB transport and Airy matching]\label{prop:largeell_wkb_matching}
Fix $J\Subset[-\mu,\mu]$ and $K\Subset(r_+,\infty)$.  For each polarization $P\in\{-1,0,+1\}$ and each sufficiently small $h$, there exist exact channel solutions $w_{P,\pm}^{\mathrm{mid}}(r,\om;h)$ defined on the closed interval between the two turning points such that
\[
   w_{P,\pm}^{\mathrm{mid}}(r,\om;h)
   =
   \exp\!\Bigl(\pm h^{-1}\Phi_P(r,\om;h)\Bigr)
   \Bigl(a_{0,\pm}(r,\om;h)+h a_{1,\pm}(r,\om;h)\Bigr),
\]
where
\[
   \Phi_P(r,\om;h)=\int_{r_{-,P}(\om;h)}^{r}\sqrt{q_P(s,0;\om,h)}\,\dd s_*
\]
and the amplitudes satisfy uniform symbol bounds
\[
   \sup_{r\in K}\abs{\partial_r^\alpha\partial_\omega^\beta a_{j,\pm}(r,\om;h)}\le C_{\alpha\beta j}.
\]
Moreover, the transfer matrices from the horizon Frobenius basis and from the infinity Volterra basis to the basis $(w_{P,+}^{\mathrm{mid}},w_{P,-}^{\mathrm{mid}})$ are $O(\angles{\ell}^{N})$, together with one $\omega$-derivative, uniformly for $\om\in J$.
\end{proposition}

\begin{proposition}[Uniform cut-off resolvent bounds on compact radial sets]\label{prop:largeell_uniform_resolvent}
Fix $J\Subset[-\mu,\mu]$ and $K\Subset(r_+,\infty)$.  There exist integers $M_0,M_1\ge0$ such that for every polarization $P\in\{-1,0,+1\}$,
\[
   \sup_{r,r'\in K}\abs{\partial_\omega^j\cG_{\ell,P}(\om;r,r')}
   \le
   C_{K,J,j}\angles{\ell}^{M_j},
   \qquad j=0,1,
\]
uniformly for $\om\in J$ and $\ell\ge1$.  The same bound holds for the jump $\disc\,\cG_{\ell,P}(\om;r,r')$ across the massive cut.
\end{proposition}

\begin{proposition}[Uniform reconstruction of the physical coefficients]\label{prop:largeell_reconstruction}
For every compact $K\Subset(r_+,\infty)$ there exists an integer $N_{\mathrm{rec}}\ge0$ such that the physical branch-cut coefficients satisfy
\[
   \sup_{r\in K}\abs{a^{\bc}_{\ell m,P}(t,r)}
   \le
   C_K\angles{\ell}^{N_{\mathrm{rec}}}
   \sum_{Q\in\{-1,0,+1\}}\sum_{j\le1}
   \sup_{r\in K}\abs{\partial_r^j v_Q^{\bc}(t,r)},
\]
and the same estimate holds for the residue projectors acting on compactly supported initial data.
\end{proposition}

\begin{theorem}[Uniform angular summability of branch-cut kernels]\label{thm:uniform_angular}
For every compact $K\Subset(r_+,\infty)$ there exist an integer $N_0\ge0$ and a constant $C_K$ such that, for every admissible $(\ell,m,P)$,
\begin{equation}\label{eq:uniform_angular_kernel_bound}
   \sup_{r,r'\in K}\abs{u_{\ell m,P}^{\bc}(t,r,r')}
   \le
   C_K\angles{\ell}^{N_0}
   \begin{cases}
      t^{-(\nu_{\ell,P}+1)}, & \kap_*(t)\le1,\\[0.5ex]
      t^{-5/6}, & \kap_0(t)\ge1,
   \end{cases}
\end{equation}
with the same intermediate and very-late remainder structure established in Section~\ref{sec:threshold_to_time}, uniformly in $(\ell,m,P)$.  Moreover,
\begin{equation}\label{eq:uniform_angular_amplitudes}
   \sup_{r,r'\in K}
   \Bigl(
      \abs{A_{\ell,P}(r,r';Q)}
      +
      \abs{B_{\ell,P}(r,r';Q)}
   \Bigr)
   \le C_K\angles{\ell}^{N_0}.
\end{equation}
\end{theorem}

\begin{proof}[Proof of Theorem~\ref{thm:uniform_angular}]
Fix $K\Subset(r_+,\infty)$ and choose cutoffs $\chi,\widetilde\chi\in C_0^\infty((r_+,\infty))$ such that $\chi\equiv1$ on $K$ and $\widetilde\chi\equiv1$ on a neighborhood of $\operatorname{supp}\chi$.  Lemma~\ref{lem:largeell_barrier}, Proposition~\ref{prop:largeell_wkb_matching}, Proposition~\ref{prop:largeell_uniform_resolvent}, and Proposition~\ref{prop:largeell_reconstruction} isolate the new large-$\ell$ input.  The remaining argument assembles these ingredients with the fixed-mode threshold formulas.  The proof has five steps.

\smallskip
\noindent\textit{Step 1. Low angular momenta.}
Choose $\ell_0\ge1$ large enough that the semiclassical argument below applies for every $\ell\ge\ell_0$.  For the finitely many modes $0\le\ell<\ell_0$, including the monopole electric channel at $\ell=0$, the bound follows by taking the maximum of the fixed-mode constants in Theorems~\ref{thm:intermediate} and \ref{thm:late}.  Hence only the regime $\ell\ge\ell_0$ needs analysis.

\smallskip
\noindent\textit{Step 2. Semiclassical high-angular-momentum ellipticity on compact $r$-sets.}
Set
\[
   h=(\ell+\tfrac12)^{-1}.
\]
In the odd channel,
\begin{equation}\label{eq:uniform_angular_odd_scaled}
   \mathcal P_{h,0}(\om)
   :=
   -h^2\partial_{r_*}^2
   +\frac{h^2\ell(\ell+1)f(r)}{r^2}
   +h^2\bigl(f(r)\mu^2-\om^2\bigr).
\end{equation}
By Proposition~\ref{prop:even_polarization_form} and Lemma~\ref{lem:normal_form}, each even channel can be written on $\operatorname{supp}\widetilde\chi$ as
\begin{equation}\label{eq:uniform_angular_even_scaled}
   \mathcal P_{h,P}(\om)
   :=
   -h^2\partial_{r_*}^2
   +V_{0,P}(r)
   +hV_{1,P}(r,\om;h)
   +h^2V_{2,P}(r,\om;h),
   \qquad P=\pm1,
\end{equation}
where
\[
   V_{0,-1}(r)=\frac{f(r)\,\ell(\ell-1)}{(\ell+\tfrac12)^2r^2},
   \qquad
   V_{0,+1}(r)=\frac{f(r)\,(\ell+1)(\ell+2)}{(\ell+\tfrac12)^2r^2},
\]
and $V_{j,P}$ together with $\partial_\om V_{j,P}$ are uniformly bounded on $\operatorname{supp}\widetilde\chi$ for $j=1,2$.

Because $f(r)/r^2$ is strictly positive on the compact set $\operatorname{supp}\widetilde\chi$, there exists $c_K>0$ such that
\[
   \frac{f(r)}{r^2}\ge c_K
   \qquad\text{on }\operatorname{supp}\widetilde\chi.
\]
For $\ell\ge\ell_0$ one then has, uniformly in $\abs{\om}\le\mu$,
\begin{equation}\label{eq:uniform_angular_symbol_lb}
   \xi^2+\frac12c_K
   \le
   \Re\sigma(\mathcal P_{h,P}(\om))(r,\xi)
   \qquad
   \text{for }(r,\xi)\in T^*\operatorname{supp}\widetilde\chi,
\end{equation}
for every channel $P$.  Thus the cut-off channel operators are semiclassically elliptic on $\operatorname{supp}\widetilde\chi$.

Standard matrix-valued semiclassical elliptic regularity therefore gives a cut-off parametrix
\[
   Q_{h,P}(\om)\in\Psi_h^{-2}
\]
such that for every $N$,
\begin{equation}\label{eq:uniform_angular_parametrix}
   \chi\,R_{\ell,P}(\om\pm\ii0)\,\chi
   =
   Q_{h,P}(\om)
   +O(h^N):L^2\to H_h^N,
\end{equation}
uniformly for $\abs{\om}<\mu$, where $R_{\ell,P}(\om\pm\ii0)$ denotes the boundary-value resolvent in the channel $P$.  The absence of real cut poles and threshold resonances from Theorem~\ref{thm:mode_stability} guarantees that these boundary values are uniquely defined.  Since kernels of operators in $\Psi_h^{-2}$ on a one-dimensional manifold are $O(h^{-1})$ on compact sets, and since $\partial_\om\mathcal P_{h,P}(\om)=O(h^2)$, there exist integers $M_0,M_1$ such that
\begin{equation}\label{eq:uniform_cutoff_resolvent_kernel}
   \sup_{r,r'\in K}
   \abs{\partial_\om^j\cG_{\ell,P}(\om;r,r')}
   \le
   C_{K,j}\angles{\ell}^{M_j},
   \qquad j=0,1,
\end{equation}
uniformly for $\ell\ge\ell_0$, $\abs{\om}<\mu$, and every admissible $P$.

\smallskip
\noindent\textit{Step 3. Uniform control of the threshold matching coefficients.}
Fix a matching radius $R>\sup K$.  The map taking Cauchy data at $r=R$ to values on $K$ is governed by \eqref{eq:uniform_cutoff_resolvent_kernel}; hence it contributes at most polynomial factors in $\ell$.  At $r=R$, the small-$\kap$ outgoing basis is given by Lemma~\ref{lem:small_kappa_basis}, and the large-$\kap$ continuation is governed by Proposition~\ref{prop:large_kappa_disc}.  Because $R$ is fixed, the large-order bounds
\[
   I_\nu(z)=O\!\Bigl(\frac{(z/2)^\nu}{\Gamma(\nu+1)}\Bigr),
   \qquad
   K_\nu(z)=O\!\bigl(\Gamma(\nu)(z/2)^{-\nu}\bigr),
\]
valid for $z$ in compact subsets of $(0,\infty)$, together with the identity
\[
   W\bigl(\sqrt z\,I_\nu(z),\sqrt z\,K_\nu(z)\bigr)=1,
\]
show that after the explicit factor $\vp^{2\nu_{\ell,P}}$ is extracted, the remaining coefficients are polynomially bounded in $\nu_{\ell,P}\sim \ell$.  Likewise, the Whittaker connection coefficients are ratios of gamma functions; Stirling's formula in vertical strips shows that, after the oscillatory monodromy factors $\ee^{\pm2\pi\ii\kap}$ are removed, their dependence on $\nu_{\ell,P}$ is also at most polynomial.

Consequently, for some integer $N_{\mathrm{th}}$,
\begin{align}
   \sup_{r,r'\in K}
   \abs{\disc\,\cG_{\ell,P}(\om;r,r')}
   &\le
   C_K\angles{\ell}^{N_{\mathrm{th}}}
   \Bigl(
      \vp^{2\nu_{\ell,P}}
      +\kap\,\vp^{2\nu_{\ell,P}}
      +\vp^{2\nu_{\ell,P}+2}
   \Bigr),
   \qquad \kap\le1,
   \label{eq:uniform_small_kappa_disc}
   \\
   \sup_{r,r'\in K}
   \abs{\disc\,\cG_{\ell,P}(\om;r,r')}
   &\le
   C_K\angles{\ell}^{N_{\mathrm{th}}}
   \Bigl(
      1+\kap^{-1}+\vp
   \Bigr),
   \qquad \kap\ge1,
   \label{eq:uniform_large_kappa_disc}
\end{align}
with the same polynomial bound for one $\om$-derivative and for the coefficients $a_{\ell,P}(r,r')$, $b_{\ell,P}^\pm(r,r')$ in Propositions~\ref{prop:small_kappa_disc} and \ref{prop:large_kappa_disc}.  In particular,
\begin{equation}\label{eq:uniform_AB_bound}
   \sup_{r,r'\in K}
   \Bigl(
      \abs{A_{\ell,P}(r,r';Q)}
      +
      \abs{B_{\ell,P}(r,r';Q)}
   \Bigr)
   \le C_K\angles{\ell}^{N_{\mathrm{th}}}.
\end{equation}

\smallskip
\noindent\textit{Step 4. Uniform oscillatory inversion.}
The proofs of Theorems~\ref{thm:intermediate} and \ref{thm:late} use only the threshold formulas, one $\om$-derivative of the amplitudes, and the oscillatory estimates of Lemma~\ref{lem:endpoint_integral} and Proposition~\ref{prop:saddle_estimate}.  By \eqref{eq:uniform_small_kappa_disc}, \eqref{eq:uniform_large_kappa_disc}, and \eqref{eq:uniform_AB_bound}, every constant in those arguments is polynomially bounded in $\ell$.  Repeating the same endpoint and saddle calculations therefore gives
\[
   \sup_{r,r'\in K}\abs{u_{\ell m,P}^{\bc}(t,r,r')}
   \le
   C_K\angles{\ell}^{N_0}
   \begin{cases}
      t^{-(\nu_{\ell,P}+1)}, & \kap_*(t)\le1,\\[0.5ex]
      t^{-5/6}, & \kap_0(t)\ge1,
   \end{cases}
\]
with the same remainder bounds as in Theorems~\ref{thm:intermediate} and \ref{thm:late}.  Because the reduced equations are independent of $m$, the constants are uniform in $\abs{m}\le\ell$.

\smallskip
\noindent\textit{Step 5. Monopole and reconstruction of the physical coefficients.}
The monopole electric channel at $\ell=0$ was already absorbed in Step~1.  The passage from the channel variables $(v_{-1},v_0,v_{+1})$ back to the physical even/odd coefficients uses only the constant matrices $T_\ell^{\pm1}$, the near-identity diagonalizer of Lemma~\ref{lem:normal_form}, and at most one radial derivative.  On compact $r$-sets these operators have coefficients bounded by $C\angles{\ell}$, so the same polynomial estimate holds for the physical branch-cut propagators.  Enlarging $N_0$ if necessary and combining this with \eqref{eq:uniform_AB_bound} proves \eqref{eq:uniform_angular_kernel_bound} and \eqref{eq:uniform_angular_amplitudes}.
\end{proof}

\end{document}
